% !TeX root = ../../Book3.tex
% 纲要标题：### 3.3 Artin 环
% 纲要位置：工作记录/计划纲要.md，第 1991 行。
% 纲要细目（写作范围）：
% * 降链条件与有限长度；通过简单商模证明根幂零及商层的域坐标分解
% * 交换 Artin 环的局部直积分解；满足根幂为零的指数及由极大理想恢复因子
% * 幂零根及零维 Noether 环

\section{Artin 环}
\label{b3:sec:0303}

降链条件对一般模弱于有限长度，但对交换环本身有更强的后果：它迫使极大理想只有有限多个，迫使 Jacobson 根幂零，最终使 \(R\) 作为自身的模具有有限长度。这些性质又把环分成有限个 Artin 局部环。

\subsection{降链条件与有限长度}
\label{b3:subsec:0303:01}

设 \(R\) 为交换含幺环。若非零 \(R\) 模 \(M\) 只有 \(0\) 和 \(M\) 两个子模，则称 \(M\) 为\term[jiandan-mo]{简单模}[simple module]。若 \(R\) 模 \(M\) 存在有限子模链
\[
0=M_0\subsetneq M_1\subsetneq\cdots\subsetneq M_r=M,
\]
且每个 \(M_i/M_{i-1}\) 都简单，则称 \(M\) 为\term[youxian-changdu-mo]{有限长度模}[module of finite length]。上述链称为\term[hecheng-lie]{合成列}[composition series]，各商称为合成因子。第二卷定理12.2.3，即模的 Jordan–Hölder 定理，已经证明其层数与所取链无关，称为模的\term[mo-changdu]{长度}[length of a module]，记作 \(\ell_R(M)\)。第二卷定理12.2.4进一步证明：有限长度等价于同时满足升链与降链条件。对环 \(R\) 而言，这一判据作用于 \(R\) 作为自身的模；任意 Artin 模则未必有限生成。
零模的合成列没有合成因子，长度为零。
\symidx{\(\ell_R(M)\)：模的长度}

\begin{lemma}\label{b3:lem:artin-primes-radical}
设 \(R\ne0\) 为交换含幺 Artin 环。则每个素理想都极大，极大理想只有有限多个，而且 \(J(R)\) 幂零。
\end{lemma}
\begin{proof}
若 \(\mathfrak p\) 素，则整环 \(R/\mathfrak p\) 仍满足降链条件。要证明这个商是域，只需证明其中每个非零元可逆。对这样的元素 \(a\)，主理想链
\((a)\supseteq(a^2)\supseteq\cdots\) 稳定，所以 \(a^n=a^{n+1}b\) 对某个 \(n,b\) 成立。约去非零的 \(a^n\)，得到 \(ab=1\)。该整环是域，故 \(\mathfrak p\) 极大。

若极大理想有无穷多个，从中选取一列两两不同的极大理想 \(\mathfrak m_1,\mathfrak m_2,\ldots\)，令
\(I_n=\mathfrak m_1\cdots\mathfrak m_n\)。不同极大理想之间不能有包含关系，否则较小者的极大性会迫使二者相等。因此对每个 \(i\le n\)，可以选
\(a_i\in\mathfrak m_i\setminus\mathfrak m_{n+1}\)，则
\(\prod_i a_i\in I_n\setminus\mathfrak m_{n+1}\)，因为 \(\mathfrak m_{n+1}\) 素。另一方面
\(I_{n+1}\subseteq I_n\cap\mathfrak m_{n+1}\)，故
\(I_{n+1}\subsetneq I_n\)，与降链条件矛盾。因此极大理想只有
\(\mathfrak m_1,\ldots,\mathfrak m_r\)。

记 \(J=J(R)\)。降链条件给出 \(J^n=J^{n+1}\)，其中 \(n\ge1\)，但稳定并不直接意味着 \(J^n=0\)。此时还不知道 \(J^n\) 有限生成，不能直接对它使用 Nakayama 引理。改看其零化子
\(K=\operatorname{Ann}_R(J^n)\)：证明 \(K=R\) 就等价于证明 \(J^n=0\)，而幂的稳定性给出
\[
(K:J)=\operatorname{Ann}_R(J^{n+1})=K.
\]
若 \(K\ne R\)，严格包含 \(K\) 的理想至少有 \(R\)，降链条件允许从中选一个极小者 \(K'\)。由于 \(K\) 与 \(K'\) 之间没有中间理想，商模 \(Q=K'/K\) 是非零简单模。取 \(x\in K'\setminus K\)，它的像生成 \(Q\)，也就是 \(K'=K+Rx\)，所以 \(Q\) 循环且有限生成。\cref{b3:thm:nakayama}于是给出 \(JQ\ne Q\)；否则这个非零模就会被迫为零。简单性使 \(JQ\) 只能是零或整个 \(Q\)，因此 \(JQ=0\)。这表示 \(Jx\subseteq K\)，于是 \(x\in(K:J)=K\)，与所取 \(x\notin K\) 矛盾。所以 \(K=R\)，即 \(J^n=0\)。
\end{proof}

这个证明用降链条件找到简单商模 \(K'/K\)，把 Nakayama 引理所需的有限生成性落实在这个循环模上。上面的处理可参阅\cite[Theorem 3.2, p. 16]{Matsumura1989CommutativeRingTheory}。

\begin{theorem}\label{b3:thm:artin-finite-length}
设 \(R\) 为交换含幺 Artin 环。则 \(R\) 作为自身的模具有有限长度，因而是 Noether 环；每个有限生成 \(R\) 模也具有有限长度。
\end{theorem}
\begin{proof}
零环的结论直接成立，以下设 \(R\ne0\)。沿用\cref{b3:lem:artin-primes-radical}中极大理想及其幂零交的记号，不同极大理想两两互素，中国剩余定理给出
\[
R/J\cong\prod_{i=1}^r R/\mathfrak m_i.
\]
考虑有限滤过 \(R\supseteq J\supseteq\cdots\supseteq J^n=0\)。对 \(0\le j<n\)，商层 \(L_j=J^j/J^{j+1}\) 被 \(J\) 零化，因而是 \(R/J\) 模；子模和商模又继承降链条件，所以每个 \(L_j\) 是 Artin 模。这是降链条件在商层上提供的有限性，尚未用到环的 Noether 性。

记 \(K_i=R/\mathfrak m_i\)，在 \(R/J\cong\prod_{i=1}^rK_i\) 中，令 \(e_i=(0,\ldots,1,\ldots,0)\)，其中只有第 \(i\) 个坐标为一。由 \(\sum_i e_i=1\) 及 \(e_ie_h=0\)、\(i\ne h\)，有
\[
L_j=e_1L_j\oplus\cdots\oplus e_rL_j.
\]
每个 \(e_iL_j\) 的标量作用只经过第 \(i\) 个域 \(K_i\)，其 \(R\) 子模正是 \(K_i\) 线性子空间。它作为 \(L_j\) 的子模满足降链条件，故对应向量空间也满足降链条件。这样的向量空间必有限维：若含无限线性无关序列，将其扩充为基，逐次删去其中前几个向量所得的线性生成空间就构成严格降链。有限维向量空间中，将基向量逐个加入，就得到相邻商一维的子空间链；各一维商的 \(R\) 子模只有零和自身，所以是简单模。因此各 \(e_iL_j\)，进而每个 \(L_j\)，都具有有限长度。把各商层的合成列取逆像并连接，得到 \(R\) 的有限合成列。

由第二卷定理12.2.4，有限长度推出 Noether 性。若 \(M\) 由 \(s\) 个元素生成，它是 \(R^s\) 的商；有限直和与商都保留有限长度，因此 \(M\) 也具有有限长度。
\end{proof}

\subsection{交换 Artin 环的局部直积分解}
\label{b3:subsec:0303:02}

中国剩余定理已给出一个具体模型：
\[
\mathbb Z/72\mathbb Z\cong\mathbb Z/8\mathbb Z\times\mathbb Z/9\mathbb Z.
\]
右侧两个因子都是局部环，分别对应原环的极大理想 \((\bar2)\)、\((\bar3)\)。剩余类 \(e_1=9\)、\(e_2=64\) 在这个同构下分别对应 \((1,0)\)、\((0,1)\)，所以二者和为一，乘积为零，且各自平方等于自身。任意 \(\mathbb Z/72\mathbb Z\) 模 \(M\) 自动分成 \(9M\oplus64M\)，分别保留 \(2\) 主部分与 \(3\) 主部分，不需要为模另选生成元。一般交换 Artin 环也能由各极大理想恢复这样的局部因子。

\begin{theorem}\label{b3:thm:artin-decomposition}
设 \(R\ne0\) 为交换含幺 Artin 环，极大理想为
\(\mathfrak m_1,\ldots,\mathfrak m_r\)。对任意满足 \(J(R)^n=0\) 的整数 \(n\ge1\)，自然映射给出
\[
R\cong\prod_{i=1}^r R/\mathfrak m_i^n
\cong\prod_{i=1}^r R_{\mathfrak m_i}.
\]
每个因子都是非零 Artin 局部环。相应地，\(1=e_1+\cdots+e_r\) 分解为两两正交的非零幂等元，每个 \(R\) 模都有
\(M=e_1M\oplus\cdots\oplus e_rM\)。这些局部因子由各极大理想唯一确定，分解仅差排列。
\end{theorem}
\begin{proof}
固定满足 \(J(R)^n=0\) 的整数 \(n\ge1\)。要使商映射给出直积分解，需要各 \(\mathfrak m_i^n\) 两两互素，且它们的交为零。对不同的 \(i,j\)，有
\(\mathfrak m_i+\mathfrak m_j=R\)。写 \(a+b=1\)、\(a\in\mathfrak m_i\)、\(b\in\mathfrak m_j\)，展开 \((a+b)^{2n-1}\) 后每一项至少含 \(a^n\) 或 \(b^n\)，故
\(\mathfrak m_i^n+\mathfrak m_j^n=R\)。于是
\[
\bigcap_i\mathfrak m_i^n
=\prod_i\mathfrak m_i^n
=\left(\prod_i\mathfrak m_i\right)^n=J(R)^n=0.
\]
中国剩余定理给出第一个同构。

这些商因子都是局部环，因为包含 \(\mathfrak m_i^n\) 的极大理想必包含 \(\mathfrak m_i\)：对每个 \(a\in\mathfrak m_i\)，有 \(a^n\in\mathfrak m_i^n\)，再用素性即可。因 \(\mathfrak m_i\) 本身极大，这样的极大理想只能是 \(\mathfrak m_i\)。又因 \(\mathfrak m_i^n\subseteq\mathfrak m_i\subsetneq R\)，商环非零；它继承 Artin 性，所以是非零 Artin 局部环。

在上述直积中，令 \(e_i\) 对应第 \(i\) 个坐标为 \(1\)、其余为 \(0\) 的元素。它们非零、两两正交且和为 \(1\)。对每个 \(R\) 模 \(M\)，等式 \(m=\sum_i e_i m\) 给出所述分解；若 \(\sum_i m_i=0\)、\(m_i\in e_iM\)，逐次乘 \(e_j\) 得 \(m_j=0\)，故和是直和。

局部因子还可直接从原环的极大理想恢复，这将使它们不依赖所取的 \(n\)。考虑任意有限直积分解
\[
R\cong\prod_{j=1}^s B_j,
\]
其中 \(B_j\) 是非零局部环，唯一极大理想记作 \(\mathfrak n_j\)。沿此同构识别 \(R\) 与直积。若 \(\mathfrak m\) 是 \(R\) 的极大理想，各坐标幂等元在域 \(R/\mathfrak m\) 中只能取 \(0\) 或 \(1\)，且由正交性及其和为 \(1\)，恰有一个坐标幂等元取值 \(1\)。设它对应第 \(j\) 个坐标，则其他坐标因子都在 \(\mathfrak m\) 中，而第 \(j\) 个坐标上的商必须是域。因此
\[
\mathfrak m=\mathfrak m_j
=B_1\times\cdots\times\mathfrak n_j\times\cdots\times B_s.
\]
反过来，每个 \(\mathfrak m_j\) 的商是域 \(B_j/\mathfrak n_j\)，故它们恰好穷尽 \(R\) 的极大理想。

第 \(j\) 个坐标投影将 \(R\setminus\mathfrak m_j\) 映到 \(B_j\) 的单位，因而诱导映射
\[
R_{\mathfrak m_j}\longrightarrow B_j,
\qquad a/u\longmapsto a_j u_j^{-1}.
\]
它是满射，因为坐标投影已经满射。若其值为零，则 \(a_j=0\)；第 \(j\) 个坐标幂等元 \(e_j\notin\mathfrak m_j\) 满足 \(e_j a=0\)，所以 \(a/u=0\)。映射的核为零，从而 \(R_{\mathfrak m_j}\cong B_j\)。应用于中国剩余定理所得的分解，就得到 \(R_{\mathfrak m_i}\cong R/\mathfrak m_i^n\)。对任意另一组非零局部因子，极大理想集合仍恢复同一组局部化，故因子只差排列和同构。
\end{proof}

零环没有极大理想，采用空直积为零环的约定后，它对应空的局部分解。零环上的幺性模只有零模，相应的模分解也是空直和。

\subsection{幂零根与零维 Noether 环}
\label{b3:subsec:0303:03}

本小节把“零维”用于每个素理想都极大的非零环；一般 Krull 维数的定义和基本定理留在\chapref{b3:ch:09}。零环作为单独的平凡情形处理。

\begin{theorem}\label{b3:thm:noether-dimension-zero}
设 \(R\) 为非零交换含幺环。下列条件等价：
\begin{equivalences}
\item \(R\) 为 Artin 环；
\item \(R\) 为 Noether 环，且每个素理想都极大。
\end{equivalences}
此时幂零根与 Jacobson 根相等，均为幂零理想。\(R\) 约化当且仅当它是有限个域的直积。
\end{theorem}
\begin{proof}
由\cref{b3:lem:artin-primes-radical}及\cref{b3:thm:artin-finite-length}，第一项推出第二项。

反过来，设 \(R\) 为 Noether 环且每个素理想极大。第一章的\cref{b3:prop:noether-spectrum-components}表明其极小素理想只有有限多个，且每个素理想包含一个极小素理想。现在包含在两个素理想之间只能是相等，所以整个素谱就是有限个极大理想
\(\mathfrak m_1,\ldots,\mathfrak m_r\)。由\cref{b3:thm:radical-prime-intersection}，
\[
\sqrt{(0)}=\bigcap_i\mathfrak m_i=J(R).
\]
Noether 条件使幂零根有限生成；它的每个元素幂零，故由\cref{b3:prop:finite-nil-ideal}，该理想的某个幂为零。

记 \(J=J(R)\)，中国剩余定理仍给出 \(R/J\cong\prod_i R/\mathfrak m_i\)。这里各理想 \(J^j\) 因 Noether 性而有限生成，所以商层 \(J^j/J^{j+1}\) 是有限生成 \(R\) 模。它被 \(J\) 零化，因而也是有限生成的 \(R/J\) 模。与前面从降链条件得到 Artin 商层不同，这个方向依靠有限生成性：按坐标幂等元分解以后，每个域上的向量空间分量由原来有限个生成元的像生成，故有限维。因此各商层具有有限长度；沿有限的 \(J\) 幂滤过连接合成列，得到 \(R\) 具有有限长度，故满足降链条件。

最后，若 \(R\) 约化，上述公共根为零，直接得到 \(R\cong\prod_i R/\mathfrak m_i\)。有限个域的直积没有非零幂零元，反向也成立。
\end{proof}

上述两种有限性给出不同的蕴含关系。在 \(R\) 为 Artin 环的假设下，有
\[
\left.
\begin{gathered}
\#\operatorname{Max}R<\infty,\quad J^n=0,\\
J^i/J^{i+1}\;\text{满足降链条件}
\end{gathered}
\right\}
\Longrightarrow\ell_R(R)<\infty
\Longrightarrow R\;\text{为 Noether 环}.
\]
而在 \(R\) 为 Noether 环且每个素理想都极大的假设下，有
\[
\left.
\begin{gathered}
\operatorname{Spec}R\;\text{有限},\quad J=\sqrt{(0)},\quad J^n=0,\\
J^i/J^{i+1}\;\text{有限生成}
\end{gathered}
\right\}
\Longrightarrow\ell_R(R)<\infty
\Longrightarrow R\;\text{为 Artin 环}.
\]
\ProseParagraphBreak
去掉 Noether 条件，零维便不足以推出 Artin 性。\secref{b3:sec:0102}中的无限域直积 \(R=\prod_{n\ge1}k\) 已经给出反例：它的每个素理想都极大，而使前 \(n\) 个坐标为零的理想构成严格降链。
